\chapter{Processes as Proofs, Carefully}

\section*{13.1}
Both linear logic and session typing make resource/obligation use explicit rather than freely duplicable by default.

\section*{13.2}
Caires--Pfenning's session calculus or Wadler's CP are acceptable.

\section*{13.3}
The table only records structural roles; it defines neither a translation function nor preservation/adequacy theorems.

\section*{13.4}
Both connect two components through a private intermediate interface/resource that is eliminated or internalized by composition.

\section*{13.5}
Concurrent systems can have nondeterministic interleavings, observable communication, and liveness behavior not represented by an arbitrary proof-normalization relation.

\section*{13.6}
Not established in general; at most a calculus-specific correspondence under an explicit translation.

\section*{13.7}
Define a translation and prove at least typing preservation plus an operational correspondence/adequacy property.

\section*{13.8}
A cited theorem can carry external mathematical authority for its stated calculus. The lab only verifies the classification of our local claims.

\section*{13.9}
Show: if $\ptypes{\Gamma}{\Delta}{P}$ then the translated object is derivable at the translated interface, with a precise mapping of linear endpoints.

\section*{13.10}
Show that each relevant source reduction maps to one or more target cut-elimination steps (and usually a converse/adequacy direction if equivalence is claimed).

\section*{13.11}
Use an enum-valued claim field such as `EXACT_EQUIVALENCE`, `PROVED_CORRESPONDENCE`, `STRUCTURAL_ANALOGY`, `MOTIVATING_RESEMBLANCE`, each with source and scope.

\section*{13.12}
Full credit identifies one exact theorem from the cited source and explicitly says why its syntax/semantics differ from PROTO-0.
